\documentclass[10pt,a4paper]{amsart}
\usepackage[margin=19mm]{geometry}
\usepackage{amsmath,amssymb,mathtools}
\usepackage[scr=rsfs,cal=euler]{mathalfa}
\usepackage{xfrac}
\usepackage[unicode,hidelinks,bookmarks=false]{hyperref}
\newcommand{\ie}{i.e.\ }
\newcommand{\cf}{cf.\ }
\newcommand{\resp}{respectively\ }
\newcommand{\Subsection}{\subsection}
\providecommand{\nobreakdash}{\nobreak\hbox{-}}
\setlength{\emergencystretch}{2em}
\multlinegap0pt
\usepackage{bm}
\usepackage{bbm}
\usepackage{dsfont}
\usepackage{xspace}
\usepackage{cancel}
\usepackage{mathtools}
\usepackage{amsthm}
\makeatletter
\@ifundefined{theorem}{\newtheorem{theorem}{Theorem}}{}
\@ifundefined{lemma}{\newtheorem{lemma}[theorem]{Lemma}}{}
\makeatother
\newcommand{\Hull}{\mathbf{Hull}}
\newcommand{\R}{\mathbb{R}}
\newcommand{\Vol}{\mathbf{Vol}}
\title[Arranging convex bodies for maximum intersection volume: the]{Arranging convex bodies for maximum intersection volume: the sharp efficiency of centroid alignment}
\author{David Victor Feldman}
\date{Source: arXiv:2608.04516v1 (CC BY 4.0)}
\begin{document}
\maketitle

\noindent\emph{Source and licence.} Arranging convex bodies for maximum intersection volume: the sharp efficiency of centroid alignment. Version arXiv:2608.04516v1 (CC BY 4.0), distributed under the Creative Commons Attribution 4.0 International licence, as recorded in the arXiv licence metadata for this exact version and shown on the abstract page https://arxiv.org/abs/2608.04516v1. Full rights notice: licenses/arxiv-2608.04516-CC-BY-4.0.txt.\par\smallskip

\noindent\textit{Editorial scope.} Counted unit: the Lemma whose source label is \texttt{lem:mr} in the article (source0.tex lines 236--246, source label \texttt{lem:mr}), delivered together with the article's own complete original proof of it (source0.tex lines 248--302, one \texttt{proof} environment of 55 lines). No mathematics is written, repaired or re-derived: statement and proof are the article's own text line for line. Required context retained uncounted: none. Every definition, display and label of those lines is retained unchanged. Omitted from this package and not redistributed: 1--235, 247--247, 303--1036. Nothing inside a retained excerpt is omitted or altered: every retained body is delivered exactly as the article's lines are written, so no omission receipt of any kind accompanies this package. Citations: the retained excerpts carry no citation command at all, so no bibliography entry of the article is transcribed and no reference is redistributed. Reference adaptation: none was needed and none was made; every reference inside the delivered unit resolves inside it, so no unresolved reference or citation is produced. Printed slips kept literal: no printed word, symbol, hypothesis or number of the counted statement or of its proof was altered. Layout and class: the article's own class is \texttt{article} with packages including amsmath, amssymb, amsthm, geometry; the wrapper is the lane's pinned \texttt{amsart} editorial wrapper at 10pt on A4, which restates the theorem-like environments the retained text needs and adds the article's own macro definitions that the retained text needs (\texttt{\textbackslash Hull}, \texttt{\textbackslash R}, \texttt{\textbackslash Vol}), with extra wrapper packages (bm, bbm, dsfont, xspace, cancel, mathtools). The delivered numbering is the unit's own. Assets: no retained span contains a figure environment, a graphics-inclusion command, a PDF-page inclusion, a table or a TikZ picture; no image file is copied into this package, the package declares no asset dependency and no image file is redistributed with it. Licence: version arXiv:2608.04516v1 of this article is distributed under the Creative Commons Attribution 4.0 International licence, as recorded in the arXiv licence metadata for this exact version and shown on the abstract page https://arxiv.org/abs/2608.04516v1. A full rights notice is delivered at licenses/arxiv-2608.04516-CC-BY-4.0.txt. Characters: every retained excerpt is pure ASCII, so the delivered engine, XeLaTeX with its default Latin Modern text font, reports no missing character.
\begin{lemma}[Minkowski--Radon, with equality case]\label{lem:mr}
Let $K\subset\R^n$ be a compact convex body with non-empty interior and
centroid at the origin, and let $u$ be a unit vector.  Then
\[
h_K(u) \;\ge\; \frac{w_K(u)}{n+1}.
\]
Equality holds for a given $u$ if and only if $K$ is a {\em cone in
direction $u$}: that is, $K=\Hull\bigl(F\cup\{a\}\bigr)$ where
$F = K\cap\{x : x\cdot u = h_K(u)\}$ and $a$ is the unique point of $K$
in the opposite supporting hyperplane $\{x : x\cdot u = -h_K(-u)\}$.
\end{lemma}

\begin{proof}
Write $t = x\cdot u$, let $[t_0,t_1]$ be the range of $t$ on $K$ (so
$t_1=h_K(u)$, $t_0=-h_K(-u)$), and let $A(t)$ denote the
$(n-1)$-dimensional volume of the section $K\cap\{x\cdot u = t\}$.  By
Brunn's theorem $g:=A^{1/(n-1)}$ is concave on $(t_0,t_1)$, and in fact on
the closed interval $[t_0,t_1]$, the Brunn--Minkowski inequality applying to
the inclusion $(1-\sigma)S_a+\sigma S_b\subseteq S_{(1-\sigma)a+\sigma b}$
for all $a,b\in[t_0,t_1]$.  Moreover $g$ is continuous on $[t_0,t_1]$; at an
extreme value of $t$ this says that the sections converge to the
corresponding face.  Both facts are used at the endpoints below, and neither
is automatic there: concavity on the closed interval gives
$g(t_0)\le\lim_{t\downarrow t_0}g(t)$, and continuity promotes an
almost-everywhere identity to an identity.  After the
affine substitution carrying $[t_0,t_1]$ to $[0,1]$, the claim
$h_K(u)\ge w_K(u)/(n+1)$ becomes: the centroid of the measure
$g^{n-1}\,dt$ on $[0,1]$ lies no further right than $n/(n+1)$.

Choose $\lambda>0$ with
$\lambda^{n-1}\int_0^1 t^{\,n-1}\,dt=\int_0^1 g^{n-1}\,dt$.  Since $g$ is
concave and $g(0)\ge 0$, the ratio $g(t)/t$ is non-increasing on
$(0,1]$; hence $g^{n-1}-(\lambda t)^{n-1}$ is $\ge 0$ on an initial
interval $[0,\tau]$ and $\le 0$ on $[\tau,1]$.  Consequently
$(t-\tau)\bigl(g^{n-1}-(\lambda t)^{n-1}\bigr)\le 0$ pointwise, and since
the two profiles have equal total mass,
\[
\int_0^1 t\,g^{n-1}\,dt \;\le\; \lambda^{n-1}\!\int_0^1 t^{\,n}\,dt
\;=\;\frac{\lambda^{n-1}}{n+1}
\;=\;\frac{n}{n+1}\int_0^1 g^{n-1}\,dt ,
\]
which is the claimed centroid bound.  Equality forces
$g(t)=\lambda t$ on $[0,1]$.

It remains to identify the bodies with $g$ linear and vanishing at
$t_0$.  Since $K$ has interior, $\lambda>0$ and $\Vol_{n-1}(S_{t_1})>0$.
Fix $t\in(t_0,t_1)$ and put $s=(t-t_0)/(t_1-t_0)$.  For any $q\in S_{t_0}$,
convexity of $K$ gives
\[
  (1-s)q + s\,S_{t_1}\;\subseteq\;S_t ,
\]
and the left-hand side has $(n-1)$-volume $s^{\,n-1}\Vol_{n-1}(S_{t_1})$,
which equals $\Vol_{n-1}(S_t)$ because $g(t)=\lambda(t-t_0)$.  A compact
convex subset of a convex body of the same finite volume is the whole body,
so
\[
  S_t\;=\;(1-s)q + s\,S_{t_1}.
\]
The right-hand side determines $q$: if $q,q'\in S_{t_0}$ then
$(1-s)q+sS_{t_1}=(1-s)q'+sS_{t_1}$, and comparing centroids gives $q=q'$.
Hence $S_{t_0}$ is a single point $a$, and $S_t=(1-s)a+sS_{t_1}$ for every
$t\in(t_0,t_1)$; that is, $K=\Hull\bigl(S_{t_1}\cup\{a\}\bigr)$.
Conversely, for a cone in direction $u$ one computes $g$ linear directly,
and the centroid of
$t\,(\lambda t)^{n-1}dt/\!\int(\lambda t)^{n-1}dt$ sits exactly at
$n/(n+1)$.
\end{proof}

\end{document}
